\documentclass[10pt,a4paper]{amsart}
\usepackage[margin=19mm]{geometry}
\usepackage{amsmath,amssymb,mathtools}
\usepackage[scr=rsfs,cal=euler]{mathalfa}
\usepackage{xfrac}
\usepackage[unicode,hidelinks,bookmarks=false]{hyperref}
\newcommand{\ie}{i.e.\ }
\newcommand{\cf}{cf.\ }
\newcommand{\resp}{respectively\ }
\newcommand{\Subsection}{\subsection}
\providecommand{\nobreakdash}{\nobreak\hbox{-}}
\setlength{\emergencystretch}{2em}
\multlinegap0pt
\usepackage{amssymb}
\usepackage{mathtools}
\usepackage{xcolor}
\usepackage[shortlabels]{enumitem}
\usepackage{float}
\newtheorem{lemma}{Lemma}
\newcommand{\Z}{\mathbb{Z}}
\title{Localization and elliptic motivic relations}
\author{Peter Xu}
\date{Source: arXiv:2608.00300v1 (CC BY 4.0)}
\begin{document}
\maketitle

\noindent\emph{Source and licence.} Localization and elliptic motivic relations. Version arXiv:2608.00300v1 (CC BY 4.0), distributed by arXiv under the Creative Commons Attribution 4.0 International licence, http://creativecommons.org/licenses/by/4.0/, as recorded in the arXiv licence metadata for this exact version and shown on the abstract page https://arxiv.org/abs/2608.00300v1. Full licence notice: \texttt{licenses/arxiv-2608.00300-CC-BY-4.0.txt}.\par\smallskip

\noindent\textit{Editorial scope.} Counted unit: the article's lemma lemma (source0.tex lines 769--771). The article's complete original proof of it is delivered with it (source0.tex lines 772--820, one \texttt{proof} environment of 49 lines). No mathematics is written, repaired or re-derived: the statement and the proof are the article's own lines, in the article's own notation, and no retained line is substituted or re-flowed.  No further source-local context is needed: every cross-reference of every retained line resolves inside the two delivered spans.  Nothing was omitted from any retained span, so the package carries no omission record.  No retained line contains a citation, so the wrapper carries no bibliography and no bibliography entry.  No retained line contains a figure environment, an image-inclusion command or drawing code, so no image is delivered and the package declares no asset dependency.  Editorial formatting only, in addition: an AMS \texttt{amsart} preamble that restates the article's own declarations and macros used by the retained lines, namely the environments \texttt{lemma} on separate counters, together with the standard packages those lines need.  The retained environments print in this document as Lemma 1 in order of appearance, whereas the article prints the same items with the numbering of its own full document.  The complete native source of the article as distributed in the arXiv e-print for this exact version is bundled unchanged as \texttt{sources/206533/source0.tex}.
\begin{lemma}
    We have $2 \cdot \Psi \circ (\mathrm{id}^{\otimes 2} \otimes \eta)=0$ as maps $P^{\otimes 2} \to 1_{\mathcal{M}}(2)[2]$ (and similarly, by symmetry, for the other orderings of $\mathrm{id},\mathrm{id},\eta$ in the three tensor factors).
\end{lemma}

\begin{proof}
    Let, as usual, $\pi_{12}:T^3\longrightarrow T^2$ be the projection. Using the same adjunction trick as the previous proof to identify our map with an element of $H^2(T^2, \Z(2))$, our element 
    \[
        \Sigma:=\Psi\circ(\mathrm{id}^{\otimes2}\otimes\eta)
    \]
    corresponds to
    \[
        (\pi_{12})_*\Psi(\mathbf a:\mathbf b:\mathbf x).
    \]
    where here this the map $\Psi$ precomposed with $1(\mathbf{a})\otimes 1(\mathbf{b}) \otimes 1(\mathbf{x})$ for $\mathbf{a},\mathbf{b},\mathbf{x}$ the tautological coaugmentations of $P^{\otimes 3}_{T^3}$.

    For each pair $\mathbf a,\mathbf b$, we have
    \[
    \Psi(\mathbf a:\mathbf b:\mathbf x)
    =
    \Psi(-\mathbf b:-\mathbf a:\mathbf x-\mathbf a-\mathbf b).
    \]
    The automorphism $T^3\to T^3$ given by
    \[
    (\mathbf a,\mathbf b,\mathbf x)
    \mapsto
    (\mathbf a,\mathbf b,\mathbf x-\mathbf a-\mathbf b)
    \]
    is an automorphism over $T_2$, hence it does not change the trace along
    $q$, and the change of variable $\mathbf y:=\mathbf x-\mathbf a-\mathbf b$ yields
    \[
    \Sigma
    =
    (\pi_{12})_*\Psi(-\mathbf b:-\mathbf a:\mathbf y).
    \]
    Since $\Psi$ is alternating in its three arguments,
    \[
    (\pi_{12})_*\Psi(-\mathbf b:-\mathbf a:\mathbf y)
    =
    -(\pi_{12})_*\Psi(-\mathbf a:-\mathbf b:\mathbf y).
    \]
    By invariance under simultaneous inversion,
    \[
    \Psi(-\mathbf a:-\mathbf b:\mathbf y)
    =
    \Psi(\mathbf a:\mathbf b:-\mathbf y).
    \]
    Finally, $\mathbf y\mapsto-\mathbf y$ is another automorphism of the fibers of $q$, so
    \[
    \Sigma
    =
    -\Sigma, \mapsto 2\Sigma =0.
    \]
\end{proof}

\end{document}
